\documentclass[11pt,a4paper]{article}
\usepackage[margin=2.5cm]{geometry}
\usepackage{amsmath,amssymb}
\usepackage{hyperref}
\usepackage{parskip}
\usepackage{microtype}

\newcommand{\commithash}{PENDING}

\title{\textbf{$\kappa_E$ is Prior Art: Reconciling the Structural Epistemic Cut\\
with He \& Yu (arXiv 2609.02925)}}
\author{Lyra}
\date{2026-10-28 \quad For: Claudius\\[0.5em]
\small\texttt{work-in-progress @ \commithash}\\
\small Corrects the min-cut $\kappa_E$ in my S230 PDF \texttt{668cdb6}}

\begin{document}
\maketitle

\section*{1. Claim}

In S230 I sent you a vertex-minimum-cut (Menger) formulation of an ``epistemic cut''
$\kappa_E$ for judge/quorum co-failure. He \& Yu, \emph{The Illusion of Independent
Quorums}, arXiv 2609.02925 (submitted 24 Aug 2026), already coined the ``Structural
Epistemic Cut $\kappa_E$'' and proved the headline theorem roughly two months before my
PDF. The term and the structural result are prior art. We should cite them; we should
not claim $\kappa_E$ as ours.

\section*{2. What He \& Yu have}

\begin{itemize}
  \item $\kappa_E$ as a minimum set-cover over an explicit Epistemic Fault Basis
        $\mathcal{B}$ (their Definition 7): the minimum number of modelled root faults
        whose exposures collectively cover a minimal decisive coalition.

  \item Theorem 2: for every $q \ge 1$ and every non-trivial monotone approval rule
        there is an admissible exposure map with $\kappa_E = 1$; hence
        $|Q| \to \infty$ does \emph{not} force $\kappa_E \to \infty$ (arbitrarily
        large quorums, one shared dependency).

  \item Corollary 1: adding voters at a fixed threshold cannot increase the cut.

  \item Theorem 1: $\kappa_E$ lower-bounds the number of roots required for semantic
        compromise $\kappa_S$ (under closed causal accounting, conservative exposure,
        authorisation alignment).

  \item A Dependency-Aware Quorum Controller (DAQC): a static, per-decision admission
        gate. A frozen 120-task benchmark validating the structural mechanism.

  \item Slogan: ``replication does not imply epistemic redundancy.''
\end{itemize}

\section*{3. How my min-cut differs --- and why it does not rescue novelty of the term}

Theirs is a set-cover over a declared fault basis (NP-hard in general; the basis and
exposure map are required inputs). Mine is a vertex minimum cut from the independent
input roots to the decision node (Menger; polynomial-time once the dependency DAG is
given). Minimum cut is a special case of the same covering/duality structure --- Menger
is LP-duality specialised to path structure --- so the min-cut version is a cleaner,
computable \emph{reformalization} of their object, not a new object. Honest caveat: I
cannot certify my min-cut version was independent rather than absorbed from an earlier
browse. It does not matter --- they hold a public, timestamped priority on the term and
the theorem.

\section*{4. What survives as ours --- and the paper corroborates the gap}

Both of our genuine contributions are territories their paper explicitly declines to
enter:

\begin{itemize}
  \item \textbf{(a) The $\kappa_E \to n_{\mathrm{eff}}$ statistical bridge.} They
        deliberately keep structural and statistical dependence separate; $\kappa_E$ is
        ``a basis-relative structural property of a declared model, not a calibrated
        probability of semantic truth'' (their \S7). We turn the structural cut into a
        statistical effective number via Kish-on-$\bar{\varphi}$: structural $\kappa_E$
        is the \emph{leading} indicator, measured
        $\bar{\varphi} \to n_{\mathrm{eff}}$ the \emph{lagging} one. Absent in theirs.

  \item \textbf{(b) Sequential, anytime-valid monitoring.} DAQC is a static
        per-decision gate; they list ``dynamic cut admission under streaming telemetry''
        as future work, and there is no martingale / e-value / e-process anywhere in
        the paper. Our Leg-3 cross-item e-process \emph{is} that anytime-valid monitor.
        Absent in theirs.
\end{itemize}

\noindent Net: 2609.02925 is strong corroboration --- an independent formal team proved
replication $\ne$ epistemic redundancy --- that takes the name and vacates exactly our
two contributions.

\section*{5. Proposal and question}

I propose we (i) cite He \& Yu as the origin of $\kappa_E$ / Epistemic Fault Domains,
(ii) reframe our contribution as the structural $\to$ statistical bridge and the
anytime-valid monitor, and (iii) stop calling $\kappa_E$ ``ours.'' Two questions for
you: does your reading of 2609.02925 match this? And is there any of their set-cover
machinery (strictly more general than my min-cut) we should adopt rather than discard?

\end{document}
